\section{A uniform cutoff and a minimum exponent of three}
\label{sec:distinct-tail}

\begin{proposition}
\label{prop:distinct-tail}
For distinct primes $p,q,r$ and positive exponents ordered as
$2\leq a\leq b\leq c$, the graph $G_{p^a q^b r^c}$ is not
Laplacian integral whenever $c\geq4a^2-2a$.
\end{proposition}

\begin{proof}
Use the complement quintic $h=h_{a,b,c}$ of Section~\ref{sec:unit}.
Its first endpoint is the cubic
\begin{equation}
\label{eq:general-first-endpoint}
h(1)=Lc^3+Qc^2+Rc+S,
\end{equation}
where
\begin{align*}
L={}&a^2+b^2-1,\\
Q={}&a^3-4a^2b^2+2a^2b-a^2+2ab^2+ab-3a
       +b^3-b^2-3b+3,\\
R={}&(a-1)(b-1)(2ab+3a+3b-3),\\
S={}&(a-1)(b-1)(a+b-1)(ab+a+b-1).
\end{align*}
Here $L,R,S>0$. In fact the sharper condition $Lc+Q\geq0$
already gives $h(1)>0$. Put $T(a)=4a^2-2a$. Then
\[
Q+T(a)L=b^2(b-1)+(2a^2+a-3)b+K(a),
\]
where $K(a)=4a^4-a^3-5a^2-a+3$. For $a=2+u$, $u\geq0$,
\[
K(2+u)=4u^4+31u^3+85u^2+95u+37>0.
\]
Consequently $Lc+Q>0$ for $c\geq T(a)$. Since $h(0)<0$,
a root lies in $(0,1)$, producing a noninteger graph eigenvalue in
$(V-1,V)$ by~\eqref{eq:complement} and Lemma~\ref{lem:join}.
\end{proof}

\begin{theorem}
\label{thm:minimum-three}
For distinct primes $p,q,r$ and all $b,c\geq1$, the graph
$G_{p^3q^br^c}$ is not Laplacian integral. Consequently every
three-prime exponent vector with minimum entry at most three gives a
nonintegral graph.
\end{theorem}

\begin{proof}
Theorems~\ref{thm:unit}, \ref{thm:minimum-two} and \ref{thm:repeated}
leave $4\leq b<c$. Proposition~\ref{prop:distinct-tail} settles
$c\geq30$. At $a=3$, equation~\eqref{eq:general-first-endpoint} is
\begin{align*}
h_{3,b,c}(1)={}&(b^2+8)c^3+(b-1)(b^2-30b-12)c^2\\
&+6(b-1)(3b+2)c+4(b-1)(b+2)(2b+1).
\end{align*}
The remaining domain $4\leq b<c<30$ has exactly
$\binom{26}{2}=325$ pairs. Exhaustive integer evaluation of this
cubic gives $257$ positive and $68$ negative values, with no zero.
Table~\ref{tab:minimum-three-signs} lists every negative range;
all other pairs in this finite domain have positive values.
The complete polynomial and separate Horner sign certificates are in
\path{results/distinct-tail.json}; their construction is described in
Section~\ref{sec:verification}. Thus a positive value gives a root in $(0,1)$.

\begin{table}[ht]
\centering
\begin{tabular}{rrrrrrrrrrr}
\toprule
$b$&4&5&6&7&8&9&10&11&12&13\\
\midrule
$c_{\min}$&5&6&7&8&9&10&11&12&13&14\\
$c_{\max}$&13&15&16&17&17&16&16&15&14&14\\
\bottomrule
\end{tabular}
\caption{Every negative $h_{3,b,c}(1)$ in $4\leq b<c<30$;
each column includes all integers between its two endpoints.}
\label{tab:minimum-three-signs}
\end{table}

For the negative values, the second endpoint is positive except at
$(b,c)=(4,5)$. Indeed, for $v\geq0$,
\[
h_{3,4,6+v}(2)=104v^3+1238v^2+4018v+2344>0.
\]
For $b\geq5$, write $b=5+u$, $c=6+u+v$, with $u,v\geq0$.
Direct expansion gives
\begin{align*}
h_{3,5+u,6+u+v}(2)
={}&(5u^2+54u+153)v^3\\
&+(20u^3+262u^2+1179u+1863)v^2\\
&+(25u^4+426u^3+2642u^2+7011u+6624)v\\
&+10u^5+218u^4+1828u^3+7200u^2+12600u+6480>0.
\end{align*}
These regions exhaust $4\leq b<c$ except $(4,5)$. Whenever
$h(1)<0<h(2)$, a root lies in $(1,2)$, giving a noninteger graph
eigenvalue in $(V-2,V-1)$.

Finally
\[
h_{3,4,5}(x)=x^5-83x^4+2327x^3-24133x^2+60576x-42480.
\]
Here $h(1)=-3792$ and $h(3/2)=48825/32>0$. The root in
$(1,3/2)$ lifts to a graph eigenvalue in $(233/2,117)$, since $V=118$.
Every interval used contains no integer. The finite nonvanishing
certificate and the written infinite-range arguments exhaust all $b,c$;
Theorems~\ref{thm:unit} and \ref{thm:minimum-two} give the final assertion.
\end{proof}
